\chapter{Det fjerde liberale kunstfaget}

\begin{motto}
Det integrerande faget er ikkje tomt fordi det\\ ikkje eig stoffet sitt. Det er fullt fordi det bind det saman.
\end{motto}

Den mest nedlatande skuldinga mot design er ikkje at faget tek feil, men at det er \emph{tomt} — at det ikkje har noko eige stoff, berre lånt metode og lånt innhald frå andre. Designaren teiknar ein stol, men materiallæra høyrer ingeniøren til; han lagar eit grensesnitt, men kognisjonen høyrer psykologen til; han formgjev ein teneste, men organiseringa høyrer økonomen til. Kvar gong ein leitar etter kjernen i design, peikar han vidare på eit anna fag. Då ser det ut som om design ikkje er ein disiplin, men eit tomrom mellom disiplinar — eit ord for å setje saman det andre har funne ut. Dette kapitlet bestrider den slutninga. Det å \emph{integrere} er ikkje fråvær av eit eige stoff; det er eit eige stoff. Og det finst ei teoretisk grunngjeving for at design er det integrerande faget i den moderne kunnskapen, like substansielt som vitskapen og humaniora, berre med ein annan oppgåve.

\section{Buchanan og den nye liberale kunsten}

Argumentet kom skarpast frå Richard Buchanan, i ein tekst som no er kanonisk \autocite{buchanan1992wicked}. Buchanan tok utgangspunkt i ein gammal idé: dei liberale kunstane, dei faga ein fri person måtte meistre for å orientere seg i verda. I antikken og mellomalderen var det grammatikk, retorikk, logikk, og så aritmetikk, geometri, astronomi og musikk — kunstane som ikkje tente eit yrke, men danna mennesket. Buchanan sitt poeng er at slike integrerande kunstar ikkje er noko fast eingong for alle; dei skifter med kulturen sine behov. Og i den teknologiske kulturen, hevda han, har det vakse fram ein \emph{ny} liberal kunst, naudsynt for å orientere seg i ei verd der mest alt vi omgjev oss med, er menneskeskapt. Den nye kunsten er design.

Kva meinte han med det? Ikkje at design er ein syssel ved sida av dei andre, men at design fyller den same rolla i dag som retorikken fylte i antikken: kunsten å føre noko ut i verda på ein måte som verkar, overtydar, og tener menneske. Der retorikken gav form til ord for å nå fram til ein lyttar, gjev design form til ting, system og røynsler for å nå fram til den som skal bruke dei. Og som retorikken var design eit \emph{generelt} fag — det handlar ikkje om eitt emne, men om korleis emne av alle slag kan gjevast verksam form. Det er nett denne generaliteten som ser ut som tomheit for den som krev at eit fag skal eige eit avgrensa stoff. Men retorikken var ikkje tom fordi han kunne brukast på krig, kjærleik og lov; han var rik nett fordi han fanst på tvers. Det same gjeld design.

Buchanan kalla difor design eit fag utan eit eige emne i tradisjonell forstand, og gjorde dette til ein styrke, ikkje ei orsaking. Designaren sitt felt er ikkje eit stykke av verda, men ein måte å gripe heile verda an på: spørsmålet om kva som tener mennesket, gjeve dei materialane, teknologiane og føresetnadene som finst. Det er eit \emph{arkitektonisk} fag i tydinga at det byggjer saman — det tek innsikt frå materiallæra, kognisjonen og økonomien og gjev henne ein form som ingen av dei kunne gje åleine. Den som seier at design «berre» set saman, har sagt det same som om ein sa at arkitekten «berre» set saman murstein og statikk. Samansetjinga \emph{er} verket.

\section{Placements, ikkje kategoriar}

Her kjem Buchanan sitt mest originale grep, og det grepet er svaret på kvifor det integrerande ikkje er tomt. Eit vanleg fag arbeider med \emph{kategoriar}: faste omgrep med eit bestemt innhald, som ein bruker til å klassifisere verda. Biologen har «art» og «celle»; juristen har «forsett» og «aktløyse». Kategorien fortel deg på førehand kva eit fenomen er. Design, hevda Buchanan, arbeider ikkje slik. Det arbeider med det han kalla \emph{placements} — plasseringar: synsstader ein medvite tek inn over eit problem for å sjå det på nytt og opne for løysingar som ikkje var synlege før.

Skilnaden er avgjerande. Ein kategori \emph{stengjer}: han fortel deg kva noko er, og dermed kva det ikkje er. Ei plassering \emph{opnar}: ho er ikkje ein dom over kva noko er, men ein invitasjon til å sjå det som om det var dette — og så sjå kva som då blir mogleg. Designaren som ser ein sjukehuskorridor snart som ein logistikk, snart som ei oppleving, snart som eit møte mellom redde menneske, har ikkje funne ut kva korridoren «eigentleg» er; han har bruka tre plasseringar til å hente fram tre sett av løysingar. Plasseringa er produktiv: ho lagar nye moglegheiter i staden for å sortere gamle. Det er dette som gjer designtenkinga generativ i staden for klassifiserande, og det er nett difor ho ikkje kan reduserast til faga ho lånar frå. Ingeniøren, psykologen og økonomen har kategoriane; designaren har kunsten å plassere problemet slik at deira kunnskap kan kome til nytte saman, mot eit menneskeleg føremål.

Og denne kunsten har eit innhald som lèt seg lære og prøve. Å plassere godt er ikkje vilkårleg; ein kan gjere det treffande eller bomma, opne fruktbart eller villleiande. Det finst betre og dårlegare plasseringar, og det finst ei opplæring i å finne dei — heile atelierpedagogikken handlar om lite anna. Den som meiner integrasjon er tomt, har forveksla «har ikkje eit fast emne» med «har ikkje noko å kunne». Design har mykje å kunne. Det er berre ikkje eit emne; det er ein kompetanse i å gje form til det ubestemte.

\section{Margolin og design studies som substans}

Om Buchanan gav den filosofiske grunngjevinga, var det Victor Margolin som synte at design òg har eit \emph{akademisk} felt med eige intellektuelt innhald \autocite{margolin2002politics}. Margolin sitt prosjekt var å etablere \emph{design studies} som eit sjølvstendig forskingsfelt, ikkje eit vedheng til kunsthistoria eller teknologistudia. Og innvendinga han måtte møte, var den same som dette kapitlet handlar om: om design lånar frå alle andre fag, kva er då igjen til design studies sjølv? Kva skulle feltet handle om, om ikkje om det andre fag alt dekkjer?

Margolin sitt svar var at design studies har eit eige objekt som ingen andre fag eig: det menneskeskapte i si fulle breidd, sett ikkje som teknisk gjenstand, ikkje som estetisk objekt, ikkje som økonomisk vare, men som det det er — den materielle og immaterielle verda menneske lagar for å leve i. Ingen disiplin tek dette som sitt: teknologistudia ser teknikken, kunsthistoria ser uttrykket, økonomien ser marknaden, sosiologien ser bruken. Design studies ser \emph{tinget i si heilskap} — korleis det vart unnfanga, forma, produsert, bruka, og kva det gjorde med liva til dei som lever med det. Det er eit genuint tverrfagleg felt, men tverrfagleg i ei sterk tyding: ikkje at det er sett saman av bitar frå andre fag, men at det har eit objekt som krev fleire fag for å bli sett, og som difor må ha sitt eige fag til å samle synet. Det integrerande er her ikkje ein mangel på fokus; det \emph{er} fokuset.

Margolin la dessutan til noko Buchanan rørte mindre ved: at dette feltet har ein etisk og politisk innsats. Tinga vi lagar, er ikkje nøytrale; dei fordeler makt, dei opnar nokre liv og stengjer andre, dei formar kva slags samfunn som blir mogleg. Eit fag som studerer det menneskeskapte, kan difor ikkje vere reint deskriptivt; det må spørje kva vi bør lage, og for kven. Difor heiter boka hans \emph{The Politics of the Artificial}. Det kunstige er aldri berre teknisk; det er alltid eit val om korleis vi vil leve. Og det er ein innsats ingen av nabofaga tek på seg i denne forma: ingeniøren spør om det verkar, økonomen om det løner seg, men spørsmålet om kva tinget gjer med menneskelivet som heilskap, fell mellom dei — og er nett det design studies tek opp. Dette er ikkje eit tomt felt. Det er eit felt med eit objekt, ein metode for å samle fleire blikk på objektet, og ein etisk forpliktelse til å spørje kva det objektet bør vere.

\section{Det integrerande som arkitektonisk logikk}

Set ein Buchanan og Margolin saman med utgangspunktet hos \textcite{simon1969sciences}, teiknar det seg eit heilt anna bilete enn tomrommet. Simon plasserte design som vitskapen om det kunstige, ved sida av — ikkje under — naturvitskapane: eit fag for det kontingente, for korleis ting kan bli. Buchanan gav dette ein arkitektonisk logikk: design er den nye liberale kunsten som bind kunnskapen saman mot menneskelege føremål, slik retorikken eingong gjorde. Margolin gav det eit akademisk objekt og ein etisk innsats. Til saman seier dei tre at design ikkje står i \emph{tomrommet} mellom faga, men i \emph{krysspunktet} — og at krysspunktet er ein stad, ikkje eit fråvær av stad.

Tenk på korleis kunnskapen vår er ordna. Vitskapen analyserer: han bryt verda ned i delar for å forstå korleis ho heng saman. Humaniora tolkar: ho spør kva delane tyder for mennesket. Men nokon må òg \emph{setje saman att} — ta det analysen og tolkinga gjev, og gjere det til noko som faktisk kan leve i verda, brukast av menneske, tene eit liv. Det er den syntetiske oppgåva, og ho er ikkje mindre intellektuell enn analyse og tolking; ho er berre annleis. Ho krev at ein held mange omsyn i hovudet samstundes, veg dei mot kvarandre når dei dreg i ulike retningar, og finn ei form der dei kan sameksistere. Dette er ikkje noko som blir til av seg sjølv når analysen er ferdig. Det er eit eige arbeid, med eigne dygder — dømmekraft, integrasjonsevne, kjensle for heilskap — og design er faget som dyrkar det arbeidet medvite. \textcite{cross2006} bygde dette ut til ein heil epistemologi: at det syntetiske, formgjevande, framtidsretta er ei eiga kunnskapskultur med eigne måtar å vite på, ikkje ein lågare versjon av vitskapen.

Då snur skuldinga seg. At design «peikar vidare» på andre fag når ein leitar etter kjernen, er ikkje prov på at det ikkje har nokon kjerne. Det er prov på kva slags kjerne det har: ein integrerande kjerne, hvis heile vesen er å hente saman. Eit fag som ikkje peika vidare, ville ikkje vere eit integrerande fag i det heile. Å klage på at design lånar frå andre, er som å klage på at brua treng to bredder — utan dei to breddene er det inga bru, men brua er likevel ikkje «berre» bredder. Ho er det som spenner over. Design er det som spenner over, og spennet er ikkje tomrom. Det er bygd.

\vignett

Men eit fag treng meir enn ei plassering i kunnskapen sitt hus; det treng eit eige stoff å gje form til. Og kva er meir designaren sitt eige enn meininga ting ber? Det er emnet for neste kapittel.

\vidarelesing{\fullcite{buchanan1992wicked}; \fullcite{margolin2002politics}; \fullcite{simon1969sciences}; \fullcite{cross2006}.}
